% TOP_PROBLEM: {"problemId": "problem.rl-versus-l", "canonicalTitle": "RL versus L", "releaseRank": 253, "primaryDomain": "theoretical_computer_science", "primaryDomainLabel": "Theoretical computer science", "displayStatus": "open", "releaseStatus": "open"}
% !TeX program = lualatex
% Definition and sources only; this file is not an LLM solution attempt.
\documentclass[11pt]{article}
\usepackage[margin=25mm]{geometry}
\usepackage{fontspec}
\setmainfont{Latin Modern Roman}
\usepackage{amsmath,amssymb,mathtools}
\usepackage{unicode-math}
\setmathfont{Latin Modern Math}
\usepackage{xurl,hyperref}
\hypersetup{colorlinks=true,urlcolor=blue}
\setcounter{secnumdepth}{0}
\setlength{\parindent}{0pt}
\setlength{\parskip}{5pt}
\setlength{\emergencystretch}{3em}
\begin{document}
\section*{\texorpdfstring{RL versus L}{RL versus L}}\label{253}
\subsection{Definitions and mathematical statement}
On length-$N$ binary inputs, $\mathrm L$ uses a deterministic machine with read-only input and $O(\log(N+2))$ writable work bits. In the selected $\mathrm{RL}$ model, random bits are unbiased and read once, and the machine halts for every input and every random tape. It must satisfy
\[
 x\in L\Rightarrow\Pr[\mathrm{accept}]\ge2/3,
 \qquad x\notin L\Rightarrow\Pr[\mathrm{accept}]=0.
\]
The question is whether every language with such a randomized logarithmic-space algorithm also has a deterministic logarithmic-space algorithm. Retaining old random bits requires counted workspace. The exact all-random-tapes halting convention is retained, rather than replaced by almost-sure termination or an expected-space bound.
\par\smallskip\textit{Formulation and concept references:} \href{https://eccc.weizmann.ac.il/report/2022/121/revision/1/download/}{[S1]}.

\subsection{Short English statement}
Can one-sided-error randomness always be removed from a logarithmic-memory decision algorithm under the stated halting and random-bit conventions?
\subsection{Sources}
\begin{itemize}
\item[S1]William M. Hoza, Recent Progress on Derandomizing Space-Bounded Computation; ECCC TR22-121 Revision1.
\url{https://eccc.weizmann.ac.il/report/2022/121/revision/1/download/}.
\textit{Formulation source.}
\end{itemize}
\end{document}
